%% kapa: exam=1 round=36 year=2025 period=1 form=A subject=civil_law subject_ko=민법 seq=40 number=40
%% source: Q-Net 감정평가사 기출문제 (cst00309) artlSeq=5231525 file=raw/1008/5231525/2025년 제36회 감정평가사 1차시험 1교시 문제지 원본.pdf
%% conversion: kordoc 4.24.51 -> tools/gen1.py (auto); review=proofread-sonnet; flags=none
\begin{question}{40}
乙은 甲에게 1억 원을 빌려주면서 甲의 X토지와 丙의 Y건물에 공동저당권을 설정받았다. 이후 甲은 X토지를 丁에게 매도하고 소유권을 이전해 주었다. 채무를 변제받지 못한 乙이 저당권을 실행하고자 한다. 이에 관한 설명으로 옳은 것은? (다툼이 있으면 판례에 따름)
\begin{choices}
  \item 乙이 Y건물의 경매에서 먼저 배당을 받을 수 있는 금액은 X토지와 Y건물의 경매대가에 비례하여 丙이 분담하는 채권액에 한한다.
  \item 乙이 X토지의 경매에서 먼저 배당을 받은 경우, X토지의 후순위저당권자는 Y건물에 대한 乙의 저당권을 대위행사할 수 있다.
  \item Y건물의 경매를 막기 위해 甲의 채무를 모두 변제한 丙은 X토지에 대한 乙의 저당권을 대위행사할 수 없다.
  \item X토지와 Y건물이 경매되어 동시에 배당이 되는 경우, 乙은 X토지의 경매대가에서 먼저 배당받고 변제받지 못한 부분에 한하여 Y건물의 경매대가에서 배당받는다.
  \item Y건물이 먼저 경매되고 X토지가 경매된 경우, 丁은 자신이 X토지에 지출한 유익비가 있어도 乙의 저당권을 대위하는 丙보다 우선하여 상환받을 수 없다.
\end{choices}
\end{question}
\answer{4}
